\chapter{Preposiciones básicas}
\section{Objetivos}
\begin{itemize}[leftmargin=*]
	\item Usar preposiciones de lugar y tiempo más frecuentes en nivel A1.
	\item Integrarlas en frases con orden V2 y negación básica.
\end{itemize}

\section{Contenidos clave}
\begin{itemize}[leftmargin=*]
	\item Lugar: \textit{in, op, onder, boven, naast, tussen, voor, achter, bij}.
	\item Tiempo: \textit{om, in, op, sinds, tot, van \ldots{} tot \ldots{}} (formatos simples).
	\item Combinación con artículos contractos comunes (exposición): \textit{in de = in de, op het = op het / op \textquotesingle{}t}.
\end{itemize}

\section{Ruta del capítulo}
Seis secciones con definición, usos de lugar y tiempo, y prevención de errores.

\section{Qué son las preposiciones}
\textbf{Función.} Conectan un sustantivo o pronombre con el resto de la frase para indicar lugar, tiempo o relación.

\textbf{Patrones A1}
\begin{itemize}[leftmargin=*]
  \item Lugar: \textit{in, op, onder, naast, voor, achter, bij}. 
  \item Tiempo: \textit{om, in, op}. 
\end{itemize}

\textbf{Mini práctica}
\begin{itemize}[leftmargin=*]
  \item Identifica la preposición en cinco frases cortas y di si es de lugar o tiempo.
  \item Completa huecos con la preposición correcta (3 oraciones).
\end{itemize}

\section{Preposiciones de lugar (\textit{in, op, onder, naast})}
\textbf{Usos básicos.}
\begin{itemize}[leftmargin=*]
  \item \textit{in} = dentro: \textit{in de tas}, \textit{in de klas}.
  \item \textit{op} = sobre/superficie: \textit{op tafel}, \textit{op het bureau}.
  \item \textit{onder} = debajo: \textit{onder de stoel}.
  \item \textit{naast} = al lado de: \textit{naast het raam}.
\end{itemize}

\textbf{Ejemplos}
\begin{itemize}[leftmargin=*]
  \item \textit{De laptop ligt op de tafel}.
  \item \textit{De sleutels zijn in mijn jas}.
  \item \textit{De kat slaapt onder het bed}.
\end{itemize}

\textbf{Mini práctica}
\begin{itemize}[leftmargin=*]
  \item Describe la posición de tres objetos en tu escritorio.
  \item Corrige una frase con preposición errónea (crea tu propio error).
\end{itemize}

\section{Preposiciones de tiempo (\textit{in, om, op})}
\textbf{Usos rápidos.}
\begin{itemize}[leftmargin=*]
  \item \textit{om} + hora exacta: \textit{om drie uur}. 
  \item \textit{op} + días/fechas: \textit{op maandag}, \textit{op 5 mei}. 
  \item \textit{in} + meses/estaciones/partes del día: \textit{in juli}, \textit{in de winter}, \textit{in de ochtend}. 
\end{itemize}

\textbf{Ejemplos}
\begin{itemize}[leftmargin=*]
  \item \textit{De les begint om acht uur}. 
  \item \textit{We zien elkaar op vrijdag}. 
  \item \textit{Ik ga op 10 juni op vakantie}. 
\end{itemize}

\textbf{Mini práctica}
\begin{itemize}[leftmargin=*]
  \item Escribe tres frases: una con \textit{om}, otra con \textit{op} y otra con \textit{in}.
  \item Pasa de fecha numérica a frase completa con preposición.
\end{itemize}

\section{Uso con lugares comunes}
\textbf{Combinaciones frecuentes.}
\begin{itemize}[leftmargin=*]
  \item \textit{op school}, \textit{op kantoor}, \textit{op het station}.
  \item \textit{in de supermarkt}, \textit{in de stad}, \textit{in de klas}.
  \item \textit{bij de bushalte}, \textit{bij de dokter}, \textit{bij vrienden}.
\end{itemize}

\textbf{Ejemplos}
\begin{itemize}[leftmargin=*]
  \item \textit{Ik ben op kantoor tot vijf uur}.
  \item \textit{Hij wacht bij de bushalte}.
  \item \textit{Wij wonen in de stad}.
\end{itemize}

\textbf{Mini práctica}
\begin{itemize}[leftmargin=*]
  \item Escribe cuatro frases con lugar y preposición fija correcta.
  \item Cambia \textit{in de stad} a otro lugar manteniendo la estructura.
\end{itemize}

\section{Uso con tiempo y horarios}
\textbf{Estructuras modelo.}
\begin{itemize}[leftmargin=*]
  \item \textit{Ik werk van negen tot vijf}. 
  \item \textit{De winkel is open van acht uur tot zes uur}. 
  \item \textit{In het weekend slaap ik tot laat}. 
\end{itemize}

\textbf{Contrastes}
\begin{itemize}[leftmargin=*]
  \item Intervalos: \textit{van \\_ tot \\_}. 
  \item Momento puntual: \textit{om \\_}. 
  \item Periodo amplio: \textit{in de zomer}, \textit{in 2026}. 
\end{itemize}

\textbf{Mini práctica}
\begin{itemize}[leftmargin=*]
  \item Escribe tu horario laboral con \textit{van \\_ tot \\_}.
  \item Cambia una frase de momento puntual a intervalo (y viceversa).
\end{itemize}

\section{Errores comunes}
\begin{itemize}[leftmargin=*]
  \item Confundir \textit{op} y \textit{in} con lugares: *\textit{op de supermarkt} $\rightarrow$ \textit{in de supermarkt}. 
  \item Olvidar \textit{om} con horas: *\textit{De les begint acht uur} $\rightarrow$ \textit{De les begint om acht uur}. 
  \item Usar \textit{in} con días: *\textit{in maandag} $\rightarrow$ \textit{op maandag}. 
  \item Duplicar preposiciones: *\textit{op in de tafel}. Usa solo una: \textit{op de tafel}. 
  \item Orden V2: coloca el verbo en segunda posición incluso con complemento inicial. 
\end{itemize}


\section{Pronunciación}
\begin{itemize}[leftmargin=*]
	\item \textit{bij} suena /bɛi̯/; evita /bi/.
	\item \textit{voor} mantiene vocal larga /voːr/; no acortes.
\end{itemize}

\section{Errores comunes}
\begin{itemize}[leftmargin=*]
	\item Traducir literalmente \textit{en} por \textit{in} cuando es \textit{op} (\textit{op straat}, \textit{op kantoor}).
	\item Olvidar \textit{om} con horas: \textit{om acht uur}.
	\item Duplicar preposición en preguntas (\textit{Waar ga je naar naartoe?}); usa solo una.
\end{itemize}

\section{Práctica guiada}
\begin{itemize}[leftmargin=*]
	\item Empareja imágenes con preposición correcta.
	\item Completa horarios con \textit{om/in/op} y rangos \textit{van ... tot ...}.
	\item Corrige frases con preposición incorrecta.
\end{itemize}

\section{Ejercicios de producción}
\begin{itemize}[leftmargin=*]
	\item Describe tu habitación usando al menos 8 preposiciones de lugar.
	\item Escribe tu agenda de un día con preposiciones de tiempo adecuadas.
\end{itemize}

\section{Mini repaso}
\begin{itemize}[leftmargin=*]
	\item Eliges \textit{in/op/bij} correctamente según contexto.
	\item Usas \textit{om} con horas y \textit{op} con fechas/días.
	\item Evitas duplicar preposiciones en preguntas y respuestas.
\end{itemize}

\section{Resultado esperado}
\begin{itemize}[leftmargin=*]
	\item Al terminar puedes ubicar objetos/personas y hablar de horarios simples usando preposiciones correctas y pronunciación natural.
\end{itemize}

\section{Práctica sugerida}
\begin{itemize}[leftmargin=*]
	\item Describir posiciones en imágenes y planos.
	\item Completar horarios usando preposiciones temporales.
	\item Mini historias con marcadores de lugar/tiempo.
\end{itemize}

\section{Microevaluación}
\begin{itemize}[leftmargin=*]
	\item 12 frases mezclando preposiciones de lugar y tiempo correctamente.
\end{itemize}
